% M. Kac, "On the notion of recurrence in discrete stochastic processes", Bull. Amer. Math. Soc. 53 (1947), first page, reset in Mills 8A.
\documentclass[11pt]{article}
\usepackage{bulletin1947}
\begin{document}
\setcounter{page}{1002}
\received{Received by the editors March 3, 1947.}
\articletitle{ON THE NOTION OF RECURRENCE IN DISCRETE\\ STOCHASTIC PROCESSES}{m. kac\footnote{John Simon Guggenheim Memorial Fellow.}}
\head{1. Introduction.} It is the purpose of this note to discuss ``statistical'' versions of the ``Wiederkehrsatz'' and the Poincar\'e cycle and to relate these versions to the ones encountered in dynamical considerations.\footnote{That the theory of stationary stochastic processes is mathematically equivalent with an ``ergodic'' theory (to which one is also led by dynamical considerations) was clearly recognized by Doob in 1934. See J. L. Doob, \emph{Stochastic processes and statistics}, Proc.\ Nat.\ Acad.\ Sci.\ U.S.A. vol.~20 (1934) pp.~376--379. The analogies discussed in the present paper are but particular cases of Doob's general equivalence principle. However, since the motivations underlying the statistical and the dynamical points of view are of a somewhat different physical character it seemed desirable to treat both cases separately.} Although the content of the note is entirely elementary and in part known it is hoped that it will help elucidate some of the basic notions of statistical physics.

\medskip
\head{2. Recurrence and mean recurrence time in a class of discrete stochastic processes.} Let $x_1, x_2, \cdots$ be a sequence of random variables each capable of assuming the values $a_1, a_2, \cdots$. We shall say that the sequence $x_1, x_2, \cdots$ is a stationary process if: (a) for each $j$ the probability
\[
\text{Prob.\ } \{x_n = a_j\}
\]
is independent of $n$; (b) for each set of values $a_{s_1}, a_{s_2}, \cdots, a_{s_r}$ the probability
\[
\text{Prob.\ } \{x_{k_1} = a_{s_1},\ x_{k_2} = a_{s_2}, \cdots, x_{k_r} = a_{s_r}\}
\]
depends only on the differences $|k_i - k_j|$.

Let
\[
\text{Prob.\ } \{x_n = a_j\} = W_1(a_j) = W(a_j)
\]
and
\[
\text{Prob.\ } \{x_1 = a_{s_1},\ x_2 = a_{s_2}, \cdots, x_r = a_{s_r}\} = W_r(a_{s_1}, a_{s_2}, \cdots, a_{s_r}).
\]
We shall use the notation $\bar a_k$ to denote any $a_j$ different from $a_k$. For instance,
\[
\text{Prob.\ } \{x_1 = a_1,\ x_2 \neq a_2\} = W_2(a_1, \bar a_2),
\]
\end{document}
